% --- Archon physics formalization source begin ---
% archon:physics
% archon:covers PhyXMiniProblems/problem_phyx_mini_0000.lean
% archon:source-report reports/phyx_mini/problem_phyx_mini_0000.source.json

\chapter{Physics problem phyx\_mini\_0000}
\label{ch:PhyXMiniProblems_problem_phyx_mini_0000}

\paragraph{Problem source.}
\#\# Physical scenario

Figure shows a refracted light beam in linseed oil making an angle of \textbackslash{}( \textbackslash{}phi = 20.0\textasciicircum{}\{\textbackslash{}circ\} \textbackslash{}) with the normal line \textbackslash{}( NN' \textbackslash{}). The index of refraction of linseed oil is 1.48.

\#\# Question

Determine the angles\textbackslash{}( \textbackslash{}theta' \textbackslash{}).

\#\# Answer choices

- A: A: \textbackslash{}( 28.5\textasciicircum{}\{\textbackslash{}circ\} \textbackslash{})

- B: B: \textbackslash{}(  30.4\textasciicircum{}\{\textbackslash{}circ\} \textbackslash{})

- C: C: \textbackslash{}( 22.3\textasciicircum{}\{\textbackslash{}circ\} \textbackslash{})

- D: D: \textbackslash{}( 31.1\textasciicircum{}\{\textbackslash{}circ\} \textbackslash{})

\#\# Figure caption (auxiliary; use the image as primary evidence)

The image shows a diagram illustrating the refraction of light through different media: air, linseed oil, and water. 

- **Layers**: 
  - The top layer is labeled "Air."
  - The middle layer is labeled "Linseed oil."
  - The bottom layer is labeled "Water."

- **Arrows**: A blue arrow represents the path of light as it travels through the three layers, bending at each interface between the different media.

- **Angles**: 
  - The angle of incidence in the air is labeled "θ."
  - The angle of refraction as the light enters the linseed oil is labeled "ϕ₁."
  - The angle of refraction as the light enters the water is labeled "θ'."

- **Normals**: 
  - Dashed lines perpendicular to the interfaces indicate the normals at the points of refraction.
  - Points on these lines are labeled "N" at the linseed oil interface and "N'" at the water interface.

This diagram is typically used to demonstrate Snell's Law, which describes how light refracts as it passes through different media.

\#\# Dataset metadata

Geometrical Optics; Physical Model Grounding Reasoning, Spatial Relation Reasoning

\paragraph{Recorded answer/context.}
C: C: \textbackslash{}( 22.3\textasciicircum{}\{\textbackslash{}circ\} \textbackslash{})

\paragraph{Figure/image path.}
/root/proposal\_for\_physic/science-mango/phyx\_mini\_run/phyx\_data/test\_image/0.png

\paragraph{Formalization target.}
create a compiling Lean file with sorry bodies at `PhyXMiniProblems/problem\_phyx\_mini\_0000.lean`. The Lean declarations must preserve the physical quantities, dimensions or dimensional roles, figure labels, governing-law hypotheses, and final relation expressed by this problem.
Use Mathlib/Physlib names found through LeanExplore where available. If a physics API is missing, introduce faithful local abstractions rather than scalar placeholder aliases.

\begin{theorem}[Physics formalization target]
\label{thm:physics:phyx_mini_0000:target}
\lean{PhyXMiniProblems.ProblemPhyXMini0000.thetaPrime_matches_choice_C}
\uses{def:physics:phyx-mini-0000:phyxminiproblems-problemphyxmini0000-threelayerrefractionsetup, def:physics:phyx-mini-0000:phyxminiproblems-problemphyxmini0000-degrees, def:physics:phyx-mini-0000:phyxminiproblems-problemphyxmini0000-isphysicalrefractionangle, def:physics:phyx-mini-0000:phyxminiproblems-problemphyxmini0000-haspositiverefractiveindices, def:physics:phyx-mini-0000:phyxminiproblems-problemphyxmini0000-satisfiessnelllawat, def:physics:phyx-mini-0000:phyxminiproblems-problemphyxmini0000-phione, def:physics:phyx-mini-0000:phyxminiproblems-problemphyxmini0000-thetaprime, def:physics:phyx-mini-0000:phyxminiproblems-problemphyxmini0000-hasparallelinterfacegeometry, def:physics:phyx-mini-0000:phyxminiproblems-problemphyxmini0000-matchesanswertonearesttenth}
The assigned autoformalize agent should translate this physics problem into Lean declarations in the covered file, with theorem and lemma proof bodies written as `by sorry`.
\end{theorem}
\begin{proof}
This is an autoformalization task, not a proof task. Produce statements that can later be proved by the physics prover without weakening the physical model.
\end{proof}

% --- Archon declaration topology begin ---
\section{Declaration topology}
\begin{definition}[Optical Medium]
  \label{def:physics:phyx-mini-0000:phyxminiproblems-problemphyxmini0000-opticalmedium}
  \lean{PhyXMiniProblems.ProblemPhyXMini0000.OpticalMedium}
  The three optical media, in their top-to-bottom order in the figure
\end{definition}
\begin{proof}
  This is the declared model definition of Optical Medium.
\end{proof}
\begin{definition}[Refraction Interface]
  \label{def:physics:phyx-mini-0000:phyxminiproblems-problemphyxmini0000-refractioninterface}
  \lean{PhyXMiniProblems.ProblemPhyXMini0000.RefractionInterface}
  The two parallel material interfaces shown in the figure
\end{definition}
\begin{proof}
  This is the declared model definition of Refraction Interface.
\end{proof}
\begin{definition}[Normal Label]
  \label{def:physics:phyx-mini-0000:phyxminiproblems-problemphyxmini0000-normallabel}
  \lean{PhyXMiniProblems.ProblemPhyXMini0000.NormalLabel}
  Labels for the dashed normals through the two refraction points
\end{definition}
\begin{proof}
  This is the declared model definition of Normal Label.
\end{proof}
\begin{definition}[Ray Segment]
  \label{def:physics:phyx-mini-0000:phyxminiproblems-problemphyxmini0000-raysegment}
  \lean{PhyXMiniProblems.ProblemPhyXMini0000.RaySegment}
  The three segments of the single light path, one in each medium
\end{definition}
\begin{proof}
  This is the declared model definition of Ray Segment.
\end{proof}
\begin{definition}[incident Medium]
  \label{def:physics:phyx-mini-0000:phyxminiproblems-problemphyxmini0000-refractioninterface-incidentmedium}
  \lean{PhyXMiniProblems.ProblemPhyXMini0000.RefractionInterface.incidentMedium}
  \uses{def:physics:phyx-mini-0000:phyxminiproblems-problemphyxmini0000-opticalmedium, def:physics:phyx-mini-0000:phyxminiproblems-problemphyxmini0000-refractioninterface}
  This def records the incident Medium component of the problem-specific physical model
\end{definition}
\begin{proof}
  This is the declared model definition of incident Medium.
\end{proof}
\begin{definition}[transmitted Medium]
  \label{def:physics:phyx-mini-0000:phyxminiproblems-problemphyxmini0000-refractioninterface-transmittedmedium}
  \lean{PhyXMiniProblems.ProblemPhyXMini0000.RefractionInterface.transmittedMedium}
  \uses{def:physics:phyx-mini-0000:phyxminiproblems-problemphyxmini0000-opticalmedium, def:physics:phyx-mini-0000:phyxminiproblems-problemphyxmini0000-refractioninterface}
  This def records the transmitted Medium component of the problem-specific physical model
\end{definition}
\begin{proof}
  This is the declared model definition of transmitted Medium.
\end{proof}
\begin{definition}[incident Ray]
  \label{def:physics:phyx-mini-0000:phyxminiproblems-problemphyxmini0000-refractioninterface-incidentray}
  \lean{PhyXMiniProblems.ProblemPhyXMini0000.RefractionInterface.incidentRay}
  \uses{def:physics:phyx-mini-0000:phyxminiproblems-problemphyxmini0000-refractioninterface, def:physics:phyx-mini-0000:phyxminiproblems-problemphyxmini0000-raysegment}
  This def records the incident Ray component of the problem-specific physical model
\end{definition}
\begin{proof}
  This is the declared model definition of incident Ray.
\end{proof}
\begin{definition}[transmitted Ray]
  \label{def:physics:phyx-mini-0000:phyxminiproblems-problemphyxmini0000-refractioninterface-transmittedray}
  \lean{PhyXMiniProblems.ProblemPhyXMini0000.RefractionInterface.transmittedRay}
  \uses{def:physics:phyx-mini-0000:phyxminiproblems-problemphyxmini0000-refractioninterface, def:physics:phyx-mini-0000:phyxminiproblems-problemphyxmini0000-raysegment}
  This def records the transmitted Ray component of the problem-specific physical model
\end{definition}
\begin{proof}
  This is the declared model definition of transmitted Ray.
\end{proof}
\begin{definition}[normal]
  \label{def:physics:phyx-mini-0000:phyxminiproblems-problemphyxmini0000-refractioninterface-normal}
  \lean{PhyXMiniProblems.ProblemPhyXMini0000.RefractionInterface.normal}
  \uses{def:physics:phyx-mini-0000:phyxminiproblems-problemphyxmini0000-refractioninterface, def:physics:phyx-mini-0000:phyxminiproblems-problemphyxmini0000-normallabel}
  This def records the normal component of the problem-specific physical model
\end{definition}
\begin{proof}
  This is the declared model definition of normal.
\end{proof}
\begin{definition}[Three Layer Refraction Setup]
  \label{def:physics:phyx-mini-0000:phyxminiproblems-problemphyxmini0000-threelayerrefractionsetup}
  \lean{PhyXMiniProblems.ProblemPhyXMini0000.ThreeLayerRefractionSetup}
  \uses{def:physics:phyx-mini-0000:phyxminiproblems-problemphyxmini0000-opticalmedium, def:physics:phyx-mini-0000:phyxminiproblems-problemphyxmini0000-normallabel, def:physics:phyx-mini-0000:phyxminiproblems-problemphyxmini0000-raysegment}
  Physical data attached to the three-layer refraction diagram. Refractive indices are dimensionless scalar readouts. Angles are physical angles modulo a full turn, represented by Mathlib's Real.Angle
\end{definition}
\begin{proof}
  This is the declared model definition of Three Layer Refraction Setup.
\end{proof}
\begin{definition}[degrees]
  \label{def:physics:phyx-mini-0000:phyxminiproblems-problemphyxmini0000-degrees}
  \lean{PhyXMiniProblems.ProblemPhyXMini0000.degrees}
  Convert a numerical degree reading into a physical angle
\end{definition}
\begin{proof}
  This is the declared model definition of degrees.
\end{proof}
\begin{definition}[degree Readout]
  \label{def:physics:phyx-mini-0000:phyxminiproblems-problemphyxmini0000-degreereadout}
  \lean{PhyXMiniProblems.ProblemPhyXMini0000.degreeReadout}
  The representative of an acute physical angle, expressed in degrees
\end{definition}
\begin{proof}
  This is the declared model definition of degree Readout.
\end{proof}
\begin{definition}[Is Physical Refraction Angle]
  \label{def:physics:phyx-mini-0000:phyxminiproblems-problemphyxmini0000-isphysicalrefractionangle}
  \lean{PhyXMiniProblems.ProblemPhyXMini0000.IsPhysicalRefractionAngle}
  Refraction angles in this diagram lie between the ray and its normal
\end{definition}
\begin{proof}
  This is the declared model definition of Is Physical Refraction Angle.
\end{proof}
\begin{definition}[Has Positive Refractive Indices]
  \label{def:physics:phyx-mini-0000:phyxminiproblems-problemphyxmini0000-haspositiverefractiveindices}
  \lean{PhyXMiniProblems.ProblemPhyXMini0000.HasPositiveRefractiveIndices}
  \uses{def:physics:phyx-mini-0000:phyxminiproblems-problemphyxmini0000-threelayerrefractionsetup}
  The dimensionless refractive index is physically positive in every layer
\end{definition}
\begin{proof}
  This is the declared model definition of Has Positive Refractive Indices.
\end{proof}
\begin{definition}[Satisfies Snell Law At]
  \label{def:physics:phyx-mini-0000:phyxminiproblems-problemphyxmini0000-satisfiessnelllawat}
  \lean{PhyXMiniProblems.ProblemPhyXMini0000.SatisfiesSnellLawAt}
  \uses{def:physics:phyx-mini-0000:phyxminiproblems-problemphyxmini0000-refractioninterface, def:physics:phyx-mini-0000:phyxminiproblems-problemphyxmini0000-threelayerrefractionsetup}
  Snell's law at one of the two interfaces in the diagram
\end{definition}
\begin{proof}
  This is the declared model definition of Satisfies Snell Law At.
\end{proof}
\begin{definition}[theta]
  \label{def:physics:phyx-mini-0000:phyxminiproblems-problemphyxmini0000-theta}
  \lean{PhyXMiniProblems.ProblemPhyXMini0000.theta}
  \uses{def:physics:phyx-mini-0000:phyxminiproblems-problemphyxmini0000-threelayerrefractionsetup}
  The angle labelled θ in the air layer
\end{definition}
\begin{proof}
  This is the declared model definition of theta.
\end{proof}
\begin{definition}[phi One]
  \label{def:physics:phyx-mini-0000:phyxminiproblems-problemphyxmini0000-phione}
  \lean{PhyXMiniProblems.ProblemPhyXMini0000.phiOne}
  \uses{def:physics:phyx-mini-0000:phyxminiproblems-problemphyxmini0000-threelayerrefractionsetup}
  The angle labelled φ₁, measured from normal N in the oil layer
\end{definition}
\begin{proof}
  This is the declared model definition of phi One.
\end{proof}
\begin{definition}[theta Prime]
  \label{def:physics:phyx-mini-0000:phyxminiproblems-problemphyxmini0000-thetaprime}
  \lean{PhyXMiniProblems.ProblemPhyXMini0000.thetaPrime}
  \uses{def:physics:phyx-mini-0000:phyxminiproblems-problemphyxmini0000-threelayerrefractionsetup}
  The requested angle θ', measured from normal N' in the water layer
\end{definition}
\begin{proof}
  This is the declared model definition of theta Prime.
\end{proof}
\begin{definition}[Has Parallel Interface Geometry]
  \label{def:physics:phyx-mini-0000:phyxminiproblems-problemphyxmini0000-hasparallelinterfacegeometry}
  \lean{PhyXMiniProblems.ProblemPhyXMini0000.HasParallelInterfaceGeometry}
  \uses{def:physics:phyx-mini-0000:phyxminiproblems-problemphyxmini0000-threelayerrefractionsetup, def:physics:phyx-mini-0000:phyxminiproblems-problemphyxmini0000-phione}
  Parallel interfaces make the oil ray have the same angle to N and N'
\end{definition}
\begin{proof}
  This is the declared model definition of Has Parallel Interface Geometry.
\end{proof}
\begin{definition}[Answer Choice]
  \label{def:physics:phyx-mini-0000:phyxminiproblems-problemphyxmini0000-answerchoice}
  \lean{PhyXMiniProblems.ProblemPhyXMini0000.AnswerChoice}
  The four numerical answer choices printed with the problem
\end{definition}
\begin{proof}
  This is the declared model definition of Answer Choice.
\end{proof}
\begin{definition}[answer In Degrees]
  \label{def:physics:phyx-mini-0000:phyxminiproblems-problemphyxmini0000-answerindegrees}
  \lean{PhyXMiniProblems.ProblemPhyXMini0000.answerInDegrees}
  \uses{def:physics:phyx-mini-0000:phyxminiproblems-problemphyxmini0000-answerchoice}
  This def records the answer In Degrees component of the problem-specific physical model
\end{definition}
\begin{proof}
  This is the declared model definition of answer In Degrees.
\end{proof}
\begin{definition}[Matches Answer To Nearest Tenth]
  \label{def:physics:phyx-mini-0000:phyxminiproblems-problemphyxmini0000-matchesanswertonearesttenth}
  \lean{PhyXMiniProblems.ProblemPhyXMini0000.MatchesAnswerToNearestTenth}
  \uses{def:physics:phyx-mini-0000:phyxminiproblems-problemphyxmini0000-degreereadout, def:physics:phyx-mini-0000:phyxminiproblems-problemphyxmini0000-answerchoice, def:physics:phyx-mini-0000:phyxminiproblems-problemphyxmini0000-answerindegrees}
  A displayed one-decimal-place choice matches an angle when their degree readouts differ by at most half of one tenth of a degree
\end{definition}
\begin{proof}
  This is the declared model definition of Matches Answer To Nearest Tenth.
\end{proof}
\begin{lemma}[oil angle at lower interface]
  \label{lem:physics:phyx-mini-0000:phyxminiproblems-problemphyxmini0000-oil-angle-at-lower-interface}
  \lean{PhyXMiniProblems.ProblemPhyXMini0000.oil_angle_at_lower_interface}
  \uses{def:physics:phyx-mini-0000:phyxminiproblems-problemphyxmini0000-threelayerrefractionsetup, def:physics:phyx-mini-0000:phyxminiproblems-problemphyxmini0000-degrees, def:physics:phyx-mini-0000:phyxminiproblems-problemphyxmini0000-phione, def:physics:phyx-mini-0000:phyxminiproblems-problemphyxmini0000-hasparallelinterfacegeometry}
  The given 20.0° oil angle transfers from N to N' because the two interfaces, and hence their normals, are parallel
\end{lemma}
\begin{proof}
  The conclusion follows from the governing relations and figure readouts named in the statement of oil angle at lower interface.
\end{proof}
% --- Archon declaration topology end ---
% --- Archon physics formalization source end ---
